\section{Evaluation}
\label{sec:eval}

We ask four research questions:
\begin{description}
  \item[RQ1] Does the exploration loop reach every state, keep the graph
    cyclic, avoid duplicate work, and report exactly the planted violations?
  \item[RQ2] How does the perception channel's subject grounding affect audit
    accuracy, and does the error follow \cref{prop:asym}?
  \item[RQ3] Is the required-absent operation (\cref{eq:req}) necessary, or is
    an intersection-only oracle enough?
  \item[RQ4] Is the fingerprint $\phi$ stable under volatile content and still
    sensitive to real UI changes?
\end{description}

\paragraph{Subject application.}
\texttt{broken-admin} is a static six-page application (index, login,
register, dashboard, reports, admin) with four cycles and four planted
violations, browsed as \texttt{guest} (\cref{fig:app}). The policy has three
forbidden clauses (admin link, delete-user, debug panel) and one required
clause (skip-to-content link). V-01 is an Admin link rendered for every role;
the backend rejects the request, but the link itself is the exposure. V-02 is
a debug panel left in the production build. V-03 is a Delete-user control two
clicks from the seed. V-04 is a missing skip-to-content link, which is an
\emph{absence}.

\paragraph{Reproducibility.}
All numbers are produced by
\nolinkurl{docs/demo_crawl.py} (RQ1) and
\nolinkurl{docs/report/experiments/run_experiments.py} (RQ2--RQ4). The raw
output is in \nolinkurl{docs/report/data/experiments.json}. Every run is
deterministic.

\begin{figure}[t]
\centering
\begin{tikzpicture}[font=\scriptsize,
  st/.style={draw=acc,thick,rounded corners=2pt,fill=acc!6,minimum width=15mm,minimum height=6mm},
  vio/.style={st,draw=bad,fill=bad!8},
  fe/.style={-{Stealth[length=1.8mm]},thick},
  be/.style={-{Stealth[length=1.8mm]},thick,dashed,draw=bad!80}]
  \node[st]  (idx) at (0,0)      {index};
  \node[st]  (log) at (2.7,0)    {login};
  \node[st]  (reg) at (5.6,0)    {register};
  \node[vio] (das) at (2.7,-1.7) {dashboard$^{\text{V1}}$};
  \node[vio] (rep) at (0.4,-3.3) {reports$^{\text{V2}}$};
  \node[vio] (adm) at (5.0,-3.3) {admin$^{\text{V3,V4}}$};
  \draw[fe] (idx) -- (log);
  \draw[fe] (log) to[bend left=18] (reg);
  \draw[be] (reg) to[bend left=18] (log);
  \draw[fe] ([xshift=-2mm]log.south) -- ([xshift=-2mm]das.north);
  \draw[be] ([xshift=2mm]das.north) -- node[right,font=\tiny]{log out} ([xshift=2mm]log.south);
  \draw[fe] (das) to[bend right=14] (rep);
  \draw[be] (rep) to[bend right=14] (das);
  \draw[fe] (das) to[bend left=14] (adm);
  \draw[be] (adm) to[bend left=14] (das);
  \draw[be] (rep.north) to[out=100,in=200] (log.west);
  \draw[be] (adm.north east) to[out=60,in=-30] (log.south east);
\end{tikzpicture}
\caption{Exploration graph recovered for \texttt{broken-admin} (role
\texttt{guest}): 6 states, 11 edges, of which 6 are back-edges (dashed). A
tree- or DAG-based crawler would drop every dashed edge. Superscripts mark
planted violations.}
\label{fig:app}
\end{figure}

\subsection{RQ1: End-to-end exploration}
\label{sec:eval:rq1}

We run the LangGraph loop on the application with two stand-ins: an
\emph{oracle perception} that reads the planted \texttt{data-tag} annotations,
and a \emph{scripted crawler} generated from the application's
\texttt{app.json}, not a live browser. This isolates exploration and audit
logic from perception and browser behaviour. RQ1 therefore establishes the
correctness of the loop and the audit algebra, not end-to-end accuracy. RQ2
measures what happens when the oracle is removed. \Cref{tab:rq1} shows the results. The run
ends by \emph{frontier exhaustion}, not by a cap. It finds all six states and
records all 11 transitions, including 6 back-edges. The 11 executed
$(v,a)$ pairs are exactly the 11 distinct claimed pairs, so there is no
duplicate work. The four reported violations match the ground truth exactly,
with no extras. Breadth with a single browser costs 9 path replays
(21 replayed steps), a 1.9$\times$ overhead over the 11 executed actions,
which is reported rather than hidden. The full suite (545 tests passing,
4 skipped only because Neo4j and Redis were not running) checks these
properties against both loop implementations.

\begin{table}[t]
\centering\small
\caption{RQ1: end-to-end crawl of \texttt{broken-admin} (oracle perception).}
\label{tab:rq1}
\begin{tabular}{@{}lr@{}}
\toprule
Metric & Value \\
\midrule
Termination reason & frontier exhausted \\
States discovered / expected & 6 / 6 \\
Edges persisted (of which back-edges) & 11 (6) \\
$(v,a)$ executed / distinct claimed & 11 / 11 \\
Duplicate $(v,a)$ executions & 0 \\
Path replays / replayed steps & 9 / 21 \\
Actions skipped & 0 \\
Violations found / planted & 4 / 4 \\
Precision / recall & 1.00 / 1.00 \\
\bottomrule
\end{tabular}
\end{table}

\subsection{RQ2: Perception grounding and error direction}
\label{sec:eval:rq2}

We keep the negation engine fixed and swap the perception channel on all six
pages (\cref{tab:rq2}). The \emph{deterministic DOM channel} builds the
Semantic UI Map from the HTML with rule-based tagging and no annotations. It
reaches precision 0.29 and recall 0.50. The error pattern matches
\cref{prop:asym}. The channel \emph{under-reports}: it grounds the admin link
as the generic tag \texttt{admin} and the debug section as \texttt{debug}, not
as the policy subjects \texttt{admin-link} and \texttt{debug-panel}. Both
missed violations are therefore \textbf{forbidden-present false negatives}. The
same under-reporting turns the required clause into \textbf{five
required-absent false positives}: the link text ``Skip to content'' is never
matched to the subject \texttt{skip-to-content}, so the clause fires on every
page that does carry the link. The channel produces no forbidden false
positives and no required false negatives, as the proposition predicts.

The ablation \emph{+lexical grounding} normalizes hyphens and underscores in
the matcher. It removes all five false positives (precision
0.29$\to$1.00) but leaves recall at 0.50. The remaining misses are
\emph{semantic}, not lexical: deciding that a link labelled ``Admin'' is an
\texttt{admin-link}, or that a \texttt{<section>} whose heading reads
``Debug'' is a \texttt{debug-panel}, requires the kind of judgement a VLM or a
learned grounding model has to provide. This result pinpoints the research
problem that remains: \emph{subject grounding} between the vocabulary of a
policy and the vocabulary of a rendered interface.

\begin{table}[t]
\centering\small
\caption{RQ2/RQ3: audit accuracy on \texttt{broken-admin} by perception
channel and audit operator. $\text{FP}_\mathcal F/\text{FN}_\mathcal F$ and
$\text{FP}_\mathcal R/\text{FN}_\mathcal R$ split errors by clause type.}
\label{tab:rq2}
\setlength{\tabcolsep}{2.4pt}\footnotesize
\begin{tabular}{@{}lccccccc@{}}
\toprule
Configuration & P & R & F1 & $\text{FN}_\mathcal F$ & $\text{FP}_\mathcal F$ & $\text{FN}_\mathcal R$ & $\text{FP}_\mathcal R$\\
\midrule
Oracle, $\Forb\cup\Req$ & 1.00 & 1.00 & 1.00 & 0 & 0 & 0 & 0 \\
Oracle, $\Forb$ only (RQ3) & 1.00 & 0.75 & 0.86 & 0 & 0 & 1 & -- \\
DOM channel & 0.29 & 0.50 & 0.36 & 2 & 0 & 0 & 5 \\
DOM + lexical grounding & 1.00 & 0.50 & 0.67 & 2 & 0 & 0 & 0 \\
\bottomrule
\end{tabular}
\end{table}

\begin{figure}[t]
\centering
\begin{tikzpicture}
\begin{axis}[
  ybar, bar width=7pt, width=\columnwidth, height=4.6cm,
  ymin=0, ymax=1.05, ylabel={score}, ylabel style={font=\scriptsize},
  symbolic x coords={Oracle,F-only,DOM,DOM+lex},
  xtick=data, tick label style={font=\scriptsize},
  legend style={font=\scriptsize,at={(0.5,1.02)},anchor=south,legend columns=3,draw=none},
  enlarge x limits=0.18]
\addplot[fill=acc!70,draw=acc]   coordinates {(Oracle,1.00) (F-only,1.00) (DOM,0.286) (DOM+lex,1.00)};
\addplot[fill=good!60,draw=good] coordinates {(Oracle,1.00) (F-only,0.75) (DOM,0.50) (DOM+lex,0.50)};
\addplot[fill=mute!50,draw=mute] coordinates {(Oracle,1.00) (F-only,0.857) (DOM,0.364) (DOM+lex,0.667)};
\legend{Precision,Recall,F1}
\end{axis}
\end{tikzpicture}
\caption{Audit accuracy by configuration. Lexical grounding fixes precision
(required-absent false alarms) but not recall (semantic grounding of
forbidden subjects).}
\label{fig:rq2}
\end{figure}

\subsection{RQ3: Is ``required absent'' necessary?}
\label{sec:eval:rq3}
Dropping $\Req$ and auditing with the intersection $S\cap\Forb$ alone keeps
precision at 1.00 but loses V-04 (recall 0.75). No setting of the forbidden
clauses can recover it, because an absence is not an element that can
intersect anything. This is the empirical counterpart of the structural
argument in \cref{sec:model}. Accessibility obligations, consent banners, and
logout affordances are all required-absent properties.

\subsection{RQ4: Fingerprint stability and sensitivity}
\label{sec:eval:fp}
We re-render each of the six pages 50 times and inject fresh per-request noise
into every capture: \texttt{data-session}/\texttt{data-csrf} tokens, a
generated hex \texttt{id}, a \texttt{data-request-id}, a 13-digit timestamp,
and a CDP-shaped AX tree with new \texttt{nodeId}s. A naive hash of URL, DOM,
and AX produces \textbf{300 distinct states} for 6 logical pages (one per
load). $\phi$ produces exactly \textbf{6}, one per page (\cref{tab:fp}).
Normalization does not over-merge: deleting the \emph{Delete user} button from
the admin page changes its fingerprint, so a security-relevant difference
still yields a distinct state.

\begin{table}[t]
\centering\small
\caption{RQ4: distinct state ids over 50 noisy captures per page.}
\label{tab:fp}
\begin{tabular}{@{}lrr@{}}
\toprule
Page & naive hash & $\phi$ (normalized) \\
\midrule
\texttt{/index.html}            & 50 & 1 \\
\texttt{/pages/login.html}      & 50 & 1 \\
\texttt{/pages/register.html}   & 50 & 1 \\
\texttt{/pages/dashboard.html}  & 50 & 1 \\
\texttt{/pages/reports.html}    & 50 & 1 \\
\texttt{/pages/admin.html}      & 50 & 1 \\
\midrule
Total (6 logical states) & 300 & 6 \\
\multicolumn{2}{@{}l}{Real change (Delete button removed) detected?} & yes \\
\bottomrule
\end{tabular}
\end{table}

\subsection{Threats to validity}
\emph{External:} there is one small, static, hand-built application with four
planted defects. The results show the mechanism works, not that it generalises,
and \cref{sec:discussion} outlines a benchmark to address this.
\emph{Construct:} RQ1 uses an oracle perception by design. RQ2 evaluates the
deterministic channel on HTML without a live AX tree, and the VLM channel is
implemented but not evaluated here. The RQ4 noise model covers known volatile
patterns and not, for example, rotating advertisements or A/B variants.
\emph{Internal:} the lexical-grounding ablation patches the matcher's normalizer in-process
only. It measures an effect and is not a shipped change.
